\documentclass[12pt]{fphw}

\usepackage[utf8]{inputenc}
\usepackage[T1]{fontenc}
\usepackage{mathpazo}
\usepackage{graphicx}
\usepackage{booktabs}
\usepackage{tabularx}
\usepackage{amsmath}
\usepackage{xcolor}
\usepackage[section]{placeins}
\usepackage{hyperref}

\hypersetup{
    colorlinks=true,
    linkcolor=blue,
    filecolor=magenta,
    urlcolor=cyan,
}
\urlstyle{same}
\setlength{\emergencystretch}{3em}

\title{Deep Learning for NLP}
\author{sdi2200160}
\date{Spring Semester 2025-2026}
\institute{National and Kapodistrian University of Athens \\ Department of Informatics and Telecommunications}
\class{Artificial Intelligence II (YS19 M226 M262 M325 M405)}

\newcommand{\CR}{Clear Reply}
\newcommand{\AMB}{Ambivalent}
\newcommand{\CNR}{Clear Non-Reply}

\begin{document}
\maketitle

{
  \hypersetup{linkcolor=black}
  \tableofcontents
}
\newpage

\section{Abstract}

This report documents the Homework 4 system for response clarity classification on QEvasion \cite{qevasion}. The target labels remain \CR, \AMB, and \CNR, but the method changes from the single-invocation prompting of Homework 3 to the required D3-Agentic Prompting architecture. D3 decomposes classification into four specialized Qwen agents: Question Intent, Answer Content, Gap and Evasion, and Decision.

The implementation uses \texttt{Qwen/Qwen3.5-0.8B} \cite{qwen} for the required core system and compares it with a same-model single-agent prompt. This keeps the experimental contrast focused on decomposition rather than model scale. The notebook and script select the final submitted system by validation macro F1, using accuracy only as a tie-breaker.

The final local validation run is a negative result, but with an interesting change in error behavior. The same-model single-agent comparator reaches 0.3667 accuracy and 0.3062 macro F1 on a balanced 30-example validation sample, while the required D3-agentic system reaches 0.3000 accuracy and 0.2206 macro F1 after hardened label parsing. Earlier Qwen prompting runs had a very hard time detecting \AMB{}; D3 detects \AMB{}, but detects it too often.

\section{Task and Prior Assignments}

All four assignments solve the same three-way response clarity task. Homework 1 used classical vector-based classifiers, Homework 2 used encoder-only transformer fine-tuning, Homework 3 used single-invocation Qwen prompting, and Homework 4 uses a lightweight multi-agent Qwen pipeline.

\begin{table}[h]
\centering
\small
\resizebox{\textwidth}{!}{%
\begin{tabular}{lllrrl}
\toprule
Homework & System & Evaluation & Accuracy & Macro F1 & Note \\
\midrule
HW1 & TF--IDF + Logistic Regression & official test & $\sim$0.63 & $\sim$0.45 & Best classical baseline \\
HW1 & GloVe Wiki + Logistic Regression & official test & -- & $\sim$0.39 & Mean-pooled embedding baseline \\
HW1 & GloVe Twitter + Logistic Regression & official test & -- & $\sim$0.38 & Conversational embedding baseline \\
HW2 & BERT-base fine-tuning & official test & 0.6299 & 0.5561 & Best encoder-only transformer \\
HW2 & DistilBERT fine-tuning & official test & 0.5779 & 0.5066 & Smaller encoder-only transformer \\
HW2 & DeBERTa-v3 fine-tuning & official test & 0.6688 & 0.2672 & Collapsed to all \AMB{} predictions \\
HW3 & Qwen 4B few-shot prompting & balanced validation & 0.4667 & 0.4667 & Best single-invocation prompting run \\
HW3 & Qwen 0.8B few-shot prompting & balanced validation & 0.4667 & 0.4091 & Same model scale as HW4 core D3 \\
HW4 & Qwen 0.8B single-agent & balanced validation & 0.3667 & 0.3062 & Selected by macro F1 for submission \\
HW4 & Qwen 0.8B D3-agentic & balanced validation & 0.3000 & 0.2206 & Required four-agent D3 system \\
\bottomrule
\end{tabular}
}
\caption{Cross-assignment comparison. HW1 and HW2 numbers are official-test results from the earlier reports; HW3 and HW4 numbers are balanced-validation diagnostic runs and are not directly official-test comparable.}
\label{tab:previous}
\end{table}

The central recurring difficulty is \AMB. Classical methods confuse it with both endpoint classes, single-prompt Qwen systems often fail to use it, and unstable fine-tuning can collapse into predicting only \AMB. This motivates macro F1 and concrete error analysis.

\section{D3-Agentic Prompting Design}

The D3 pipeline is not repeated voting. Each agent performs a distinct linguistic subtask and passes its output to later stages.

\begin{enumerate}
    \item \textbf{Question Intent Agent}: identifies what the journalist asks for and what information a clear answer would need.
    \item \textbf{Answer Content Agent}: extracts what the answer actually says, including explicit claims, relevant details, hedges, and omissions.
    \item \textbf{Gap and Evasion Agent}: compares expected information with actual answer content and identifies missing requirements or evasion patterns.
    \item \textbf{Decision Agent}: maps the previous agents' outputs to exactly one final clarity label.
\end{enumerate}

All core agents use \texttt{Qwen/Qwen3.5-0.8B}. This is the required model in the specification and avoids a confound where one agent performs better merely because it uses a larger model. The single-agent comparator also uses the same model and taxonomy so that the comparison measures the effect of decomposition.

\section{Prompt Format}

The prompts include the original question and answer, the dataset label taxonomy, and compact JSON output constraints. The taxonomy is:
\[
\begin{aligned}
\text{\CR} &\leftarrow \text{Explicit}, \\
\text{\AMB} &\leftarrow \text{Implicit, Dodging, General, Deflection, Partial/half-answer}, \\
\text{\CNR} &\leftarrow \text{Declining to answer, Claims ignorance, Clarification}.
\end{aligned}
\]

The first three agents are instructed not to emit a final label. Only the Decision Agent may output \CR, \AMB, or \CNR. This design supports analysis: the generated \texttt{agent\_outputs.csv} files expose whether an error originates from misunderstanding the question, missing answer content, weak gap analysis, or the final mapping step.

\section{Experimental Protocol}

The implementation uses deterministic decoding, a fixed random seed, and a balanced validation sample of 10 examples per class by default. The main command is:
\begin{verbatim}
python -m artificial_intelligence.homeworks.hw20260609.sdi2200160 \
  --preset full --eval-per-label 10
\end{verbatim}

The Kaggle notebook is named:
\begin{verbatim}
sdi2200160 homework-4-d3-agentic.ipynb
\end{verbatim}
The private Kaggle notebook link will be inserted after upload.

It writes:
\begin{verbatim}
submission_best_d3_agentic_system.csv
\end{verbatim}

The full validation comparison should include:
\begin{itemize}
    \item \texttt{d3-agentic}: required four-agent system;
    \item \texttt{single-agent}: same-model direct prompt comparator;
    \item previous baselines from Table \ref{tab:previous}.
\end{itemize}

\section{Validation Results}

Table \ref{tab:hw4-results} reports the completed full run. The validation set is balanced by construction: 10 examples from each of \CR, \AMB, and \CNR. The selected system is \texttt{qwen-0.8b\_single-agent}, because it has the higher macro F1. The required D3 system is retained in the analysis because it is the assigned method and because its intermediate traces are useful for diagnosing failure modes.

\begin{table}[h]
\centering
\resizebox{\textwidth}{!}{%
\begin{tabular}{lrrrrrr}
\toprule
System & Calls & Accuracy & Macro F1 & Precision & Recall & Invalid Rate \\
\midrule
Single-agent Qwen 0.8B & 1 & 0.3667 & 0.3062 & 0.4400 & 0.3667 & 0.0000 \\
D3-agentic Qwen 0.8B & 4 & 0.3000 & 0.2206 & 0.2083 & 0.3000 & 0.0000 \\
\bottomrule
\end{tabular}
}
\caption{HW4 validation results from \texttt{runs\_hw4\_full/experiment\_summary.csv}.}
\label{tab:hw4-results}
\end{table}

\begin{table}[h]
\centering
\small
\begin{minipage}{0.45\textwidth}
\centering
\resizebox{\linewidth}{!}{%
\begin{tabular}{lrrr}
\toprule
True $\backslash$ Pred. & \CR & \AMB & \CNR \\
\midrule
\CR & 8 & 2 & 0 \\
\AMB & 10 & 0 & 0 \\
\CNR & 7 & 0 & 3 \\
\bottomrule
\end{tabular}
}
\par\smallskip\textbf{Single-agent}
\end{minipage}
\hfill
\begin{minipage}{0.45\textwidth}
\centering
\resizebox{\linewidth}{!}{%
\begin{tabular}{lrrr}
\toprule
True $\backslash$ Pred. & \CR & \AMB & \CNR \\
\midrule
\CR & 2 & 8 & 0 \\
\AMB & 3 & 7 & 0 \\
\CNR & 1 & 9 & 0 \\
\bottomrule
\end{tabular}
}
\par\smallskip\textbf{D3-agentic}
\end{minipage}
\caption{Validation confusion matrices after hardened parsing. D3 recovers more \AMB{} examples than the direct prompt, but it overuses \AMB{} and never correctly predicts \CNR{} in this run.}
\label{tab:confusion}
\end{table}

\section{Error Analysis}

The direct prompt has the better aggregate score, but its behavior is still very poor. It predicts \CR{} 25 times, \AMB{} only twice, and \CNR{} three times. Its \AMB{} recall is exactly 0.0: every true \AMB{} example is mapped to \CR. This is the same middle-class failure pattern observed in the earlier Qwen prompting assignment, where the model had trouble preserving the distinction between an incomplete answer and a clearly answered question.

The D3 system fails differently. After recovering truncated JSON label fields, it produces no invalid decisions. It predicts \AMB{} 24 times, \CR{} six times, and \CNR{} zero times. It correctly recovers seven of the ten \AMB{} examples, so the decomposition does make the model willing to use the middle class. The cost is that \AMB{} becomes too broad: all ten \CNR{} examples are missed, and most \CR{} errors are also absorbed into \AMB{}. The Gap and Evasion stage often recognizes that an answer is incomplete, but the final Decision Agent collapses many non-replies into \AMB{} rather than reserving \CNR{} for explicit refusal, ignorance, or clarification.

Failure overlap is also limited. Only about 43\% of the two systems' errors overlap, so the agentic decomposition is not merely duplicating the direct prompt's mistakes. It creates a different error surface: more interpretable and more sensitive to partialness, but less reliable as a classifier.

\section{Conclusion}

The completed system provides a reproducible D3-Agentic Prompting framework and records all intermediate agent outputs needed for interpretation. As a classifier, however, this HW4 run is substantially weaker than the earlier approaches. HW1's TF--IDF logistic regression reached about 0.45 macro F1, HW2's BERT-base fine-tuning reached 0.5561 macro F1, and HW3's best Qwen 4B few-shot prompting run reached 0.4667 macro F1 on the balanced-validation diagnostic setup. The HW4 single-agent comparator reaches only 0.3062 macro F1, and the required D3 system reaches 0.2206.

The main conclusion is therefore negative but useful: linguistic decomposition alone does not solve response-clarity classification for a very small instruction model. The interesting result is that the D3 setup reverses the earlier LLM failure mode. Previous Qwen prompting systems had a very hard time detecting \AMB{}; here, D3 detects \AMB{}, but too much. The middle label remains the central obstacle. The direct prompt effectively refuses to detect \AMB{}, while D3 overcorrects toward \AMB{} and loses \CNR{}. The final submission uses the macro-F1 selected system, but the report treats the D3 run as the core methodological result required by the assignment.

\nocite{*}
\bibliographystyle{plain}
\bibliography{refs}

\end{document}
